\section{Rodando o primeiro container e comandos de sobrevivência}

Após a instalação, é possível rodar o primeiro container com o seguinte comando:

\begin{verbatim}
docker run -d -p 8080:80 docker/welcome-to-docker
\end{verbatim}

O \textit{frontend} dessa aplicação fica acessível na porta 8080, pelo endereço \texttt{http://localhost:8080} no navegador. Após a execução, o container aparece como rodando e a imagem passa a constar nas images locais.

\subsection{Comandos básicos}

\subsubsection{Executando containers}

\begin{verbatim}
docker run -d -p 8080:80 docker/welcome-to-docker
\end{verbatim}

A flag \texttt{-d} executa o container em modo \textit{detached} (em segundo plano). A flag \texttt{-p 8080:80} mapeia a porta 8080 do host para a porta 80 do container.

\subsubsection{Listando containers}

\begin{verbatim}
docker ps          # lista containers em execução
docker ps -a       # lista todos (incluindo os parados)
\end{verbatim}

\subsubsection{Parando containers}

\begin{verbatim}
docker stop <ID>       # parar o container
\end{verbatim}

Não é necessário digitar o ID completo --- basta os caracteres suficientes para torná-lo único entre todos os containers existentes (inclusive os parados).

\subsubsection{Removendo containers}

\begin{verbatim}
docker rm <ID>         # remove um container parado
docker rm -f <ID>      # para e remove o container (força)
\end{verbatim}

O comando \texttt{docker rm} remove um container que já foi parado. A flag \texttt{-f} força a remoção mesmo que o container ainda esteja em execução --- ele é parado automaticamente antes de ser removido.

\subsubsection{Visualizando logs}

\begin{verbatim}
docker logs <ID>           # exibe todos os logs do container
docker logs -f <ID>        # segue os logs em tempo real (follow)
docker logs --tail 50 <ID> # exibe as últimas 50 linhas
\end{verbatim}

O comando \texttt{docker logs} exibe a saída padrão (\texttt{stdout}) e a saída de erro (\texttt{stderr}) de um container. A flag \texttt{-f} mantém o terminal conectado, exibindo novos logs conforme são gerados --- semelhante ao \texttt{tail -f} no Linux. Isso é útil para monitorar o comportamento de uma aplicação em tempo real.

\subsubsection{Executando comandos dentro de um container}

\begin{verbatim}
docker exec -it <ID> /bin/bash    # abre um shell interativo
docker exec <ID> ls /app          # executa um comando único
docker exec -it <ID> sh           # alternativa para imagens Alpine
\end{verbatim}

O comando \texttt{docker exec} permite executar comandos dentro de um container que já está em execução. A flag \texttt{-i} mantém o \texttt{stdin} aberto (interativo) e a flag \texttt{-t} aloca um terminal (TTY). Juntas, \texttt{-it} permitem interagir com o shell do container como se fosse um terminal remoto.

Isso é útil para depuração, inspeção de arquivos, verificação de configurações ou execução de comandos administrativos dentro do container.

\subsection{Resumo dos comandos}

\begin{center}
\begin{tabular}{|l|l|}
\hline
\textbf{Comando} & \textbf{Função} \\
\hline
\texttt{docker run} & Cria e executa um container \\
\texttt{docker ps} & Lista containers em execução \\
\texttt{docker ps -a} & Lista todos os containers \\
\texttt{docker stop <ID>} & Para um container \\
\texttt{docker rm <ID>} & Remove um container parado \\
\texttt{docker rm -f <ID>} & Para e remove um container \\
\texttt{docker logs <ID>} & Exibe logs do container \\
\texttt{docker logs -f <ID>} & Segue logs em tempo real \\
\texttt{docker exec -it <ID> sh} & Abre shell dentro do container \\
\hline
\end{tabular}
\end{center}

